% Q3: Toroidal Core — Reluctance equivalent circuit
\documentclass[border=10pt]{standalone}
\usepackage[american voltages, american currents]{circuitikz}
\usepackage{amsmath}
\renewcommand{\familydefault}{\sfdefault}

\begin{document}
\begin{circuitikz}[
    line width=0.8pt,
    every node/.style={font=\sffamily}
]

% MMF source (voltage source in reluctance analogy)
\draw (0,0) to[V, v=$NI$ (MMF)] (0,4);

% Top rail
\draw (0,4) -- (7,4);

% Core reluctance (going down)
\draw (7,4) to[R, i>^=$\Phi$] (7,2);
\node[right, xshift=6pt] at (7,3) {$\mathcal{R}_{\mathrm{core}} = \dfrac{\ell_1}{\mu_r \mu_0 A}$};

% Air gap reluctance (going down)
\draw (7,2) to[R] (7,0);
\node[right, xshift=6pt] at (7,1) {$\mathcal{R}_{\mathrm{gap}} = \dfrac{\ell_2}{\mu_0 A}$};

% Bottom return
\draw (7,0) -- (0,0);

\end{circuitikz}
\end{document}
